\section{Motivation and Research Question}
\label{sec:intro}

Data-driven weather models---GraphCast~\cite{graphcast}, GenCast, Pangu-Weather,
FourCastNet---now match or beat operational numerical weather prediction on
medium-range skill, at a small fraction of the compute. What they learned in order
to do so is unknown. A physics solver can be read term by term against the
conservation laws it discretizes; a trained graph network cannot. The opacity has
practical costs: forecasts of cyclones, heat waves, and floods feed real decisions,
and a model leaning on training-distribution correlations may fail quietly as the
climate drifts away from that distribution. There is a scientific cost too. If the
model has internalized genuine atmospheric dynamics, its internals amount to a
compressed description of those dynamics that we are currently unable to read.

A model can be highly skillful while having learned only correlational structure.
The question that separates a statistical emulator from a model of the underlying
system is causal:

\begin{quote}
Does GraphCast's internal representation encode the \emph{causal structure} of the
atmosphere, recoverable as a directed graph over interpretable features, or only
the correlational structure sufficient for next-step prediction?
\end{quote}

Accuracy does not settle this. A model at the predictability ceiling may have
organized its computation around the true dependency structure of the atmosphere,
or around shortcuts that predict equally well in distribution; observational
evaluation cannot tell the two apart, which is why interpretability claims of this
kind need interventional or otherwise causal grounding to
generalize~\cite{joshi2026}.

We approach the question along two routes. The first treats sparse-autoencoder
(SAE) features of the model's activations as variables and runs time-series causal
discovery over them. The second decomposes the model's weights into mechanisms
whose importance is defined by ablation (Sec.~\ref{sec:approach}). Neither can be
validated on GraphCast itself, where no ground-truth feature graph exists, so we
validate both on SAVAR, a synthetic spatiotemporal system whose causal graph is
known by construction (Secs.~\ref{sec:savar} and~\ref{sec:ext}).

That last sentence names a problem larger than validation. Both routes begin by
choosing variables---an aggregation of the model's state into a manageable feature
set---and on SAVAR we can pick that aggregation by trying many and keeping
whichever yields the best graph. On GraphCast we cannot, and a method that needs
the answer key to choose its own variables does not transfer. Much of the work
reported here is therefore aimed at a third question: can a candidate feature set
be ranked, and the accuracy of the graph built on it be predicted, using only
quantities computable without ground truth? Section~\ref{sec:selector} answers
yes, by crossing two independent readings of the model---one that \emph{reads}
correlations through a candidate map and one that \emph{pokes} the frozen
forecaster through it---and shows the resulting selector holds up across
overlapping features, rollout-trained forecasters, coupled channels, and the
generator's realism mechanisms, while naming the conditions under which it does
not.

\subsection{Relation to prior work}

MacMillan and Ouellette~\cite{macmillan2025} trained SAEs on GraphCast activations
and found features corresponding to recognizable physics: tropical cyclones,
atmospheric rivers, diurnal and seasonal cycles, precipitation, sea ice, geographic
position. Scaling a single tropical-cyclone feature changed the forecast in a
physically consistent way, so at least some features are causally connected to the
output. What their analysis leaves open, and what they themselves name as the open
problem, is how the features relate to one another. A single-feature intervention
establishes that a feature matters for the forecast; it says nothing about the web
of lagged influences among features, which is where any learned analogue of
atmospheric dynamics would have to live. Recovering that graph is the aim of this
project.
